\section{Related Work}
\label{sec:related}

\paragraph{Mathematical reasoning benchmarks.}
The MATH dataset~\citep{hendrycks2021measuring} evaluates competition-level mathematics across seven subjects, with problems requiring numerical or symbolic answers.
AGIEval~\citep{zhong2023agieval} includes math problems from standardized tests.
These benchmarks test \emph{problem-solving} ability but focus on high-school and undergraduate competition mathematics, where problems have unique answers and standard solution techniques.

\paragraph{Research-level mathematics.}
FrontierMath~\citep{glazer2024frontiermath} is the closest benchmark in spirit, comprising research-level mathematics problems with custom verification scripts.
However, FrontierMath spans all areas of mathematics and does not specifically target constructive reasoning; many problems still ask for unique numerical answers.
GHOSTS~\citep{zheng2024ghosts} evaluates graduate-level science tasks using rubric-based scoring, but covers mathematics only as one component of a broader scientific evaluation.

\paragraph{Formal theorem proving.}
PutnamBench~\citep{tsoukalas2024putnambench} formalizes Putnam competition problems in Lean, Isabelle, and Coq, testing whether LLMs can produce machine-verifiable proofs.
AMBER~\citep{liu2025amber} combines construction and verification in Lean~4, requiring models to both state and prove lemmas.
These benchmarks operate in the formal verification paradigm, where correctness is checked by a proof assistant.
ECAG-Bench differs by operating in the \emph{informal} mathematical setting, using CAS verification of object properties rather than formal proof checking.

\paragraph{Interactive and multi-turn evaluation.}
\citet{collins2024evaluating} propose evaluating mathematical capabilities through multi-turn interactions.
MathArena~\citep{matharena2025} uses recent competition problems with dual grading (automatic and human) to minimize contamination.
ECAG-Bench currently uses single-turn evaluation but is designed to support multi-turn interaction in future work.

\paragraph{Constructive vs.\ verificational reasoning.}
The distinction between constructive and verificational reasoning is well-established in mathematics~\citep{hartshorne1977algebraic} but has received limited attention in the LLM evaluation literature.
To our knowledge, ECAG-Bench is the first benchmark to systematically evaluate constructive reasoning in a specific mathematical domain, with property-based verification as the primary evaluation mechanism.
